% Heaviside Composite Optimization Problems
% Reconstructed with the user-provided SimplePlus Beamer theme.
% Persistent project requirements: RECONSTRUCTION_RULES.md

\documentclass[aspectratio=169,xcolor=dvipsnames]{beamer}
\usetheme{SimplePlus}

% Beamer routes math through Computer Modern Sans by default, which makes
% \mathbf render as heavy bold sans.  Restore conventional serif math while
% keeping the theme's sans-serif text.
\usefonttheme[onlymath]{serif}

\usepackage{graphicx}
\usepackage{booktabs}
\usepackage{dsfont}
\usepackage{mathtools}

% Computer Modern's italic f hooks to the right, so a following prime sits on
% top of it.  \fp inserts a small kern to separate the two.
\newcommand{\fp}{f\mkern2mu'}
\usepackage{pgfplots}
\pgfplotsset{compat=1.18}

% A calmer hierarchy than the template default: titles remain prominent
% without overwhelming dense mathematical content.
\setbeamerfont{frametitle}{size=\large,series=\bfseries,parent=structure}

% A reusable visual language for material added beyond the photographed
% lecture slides.  Original course content remains on the plain white canvas.
\definecolor{SupplementAccent}{RGB}{90,119,126}
\definecolor{SupplementTint}{RGB}{241,246,246}
\setbeamercolor{supplementary note}{fg=MediumBlack,bg=SupplementTint}
\newenvironment{supplementarycallout}[1]{%
  \begingroup
  \setbeamercolor{block title alerted}{fg=white,bg=SupplementAccent}%
  \setbeamercolor{block body alerted}{fg=black,bg=SupplementTint}%
  \begin{alertblock}{#1}%
}{%
  \end{alertblock}%
  \endgroup
}
\newenvironment{supplementarynote}{%
  \vspace{0.08cm}%
  \begin{beamercolorbox}[
    wd=\textwidth,sep=0.08cm,leftskip=0.08cm,rightskip=0.08cm
  ]{supplementary note}%
  \footnotesize
  \textcolor{SupplementAccent}{\textbf{Supplementary note.}}\enspace
}{%
  \end{beamercolorbox}%
}

% The original title page was never photographed; the speaker details below
% were supplied separately and verified against the USC Viterbi directory.
\title{Heaviside Composite Optimization Problems}
\subtitle{}
\author{Jong-Shi Pang}
\institute{Daniel J. Epstein Department of Industrial and Systems Engineering\\
  University of Southern California}
\date{Peking University, July 2026}

\begin{document}

\begin{frame}
  \titlepage
\end{frame}

\section{Chapter I: Introduction}

\begin{frame}[plain]
  \centering
  \vfill
  {\usebeamercolor[fg]{structure}\usebeamerfont{title}Chapter I\par}
  \vspace{0.34cm}
  {\usebeamercolor[fg]{structure}\usebeamerfont{subtitle}\mdseries\large
    Introduction\par}
  \vfill
\end{frame}

\begin{frame}{A bird's eye-view of post World-War-II optimization till now:}
  \small
  \begin{itemize}
    \setlength{\itemsep}{0.08cm}
    \item From linear programming to nonlinear and integer programming.
    \item From nonlinear programming to stochastic programming and
      modern-day convex programming.
    \item From differentiability to nondifferentiability.
    \item From nondifferentiability to semicontinuity.
    \item From convexity to nonconvexity.
  \end{itemize}

  \vspace{0.02cm}
  {\color{MediumBlue}\scriptsize
    Y. Cui and J. S. Pang, \textit{Modern Nonconvex Nondifferentiable
    Optimization}. SIAM, MOS--SIAM Series on Optimization
    (December 2021).\par}

  \vspace{0.12cm}
  {\color{MediumBlue}\rule{0.54\textwidth}{0.45pt}\par}
  \vspace{0.10cm}

  {\color{MediumBlue}\small
    Apart from further developments of the present, what is the future??\par}
  \vspace{0.05cm}
  {\small
    From semicontinuity to lack thereof
    $\Rightarrow$ discontinuous optimization.\par
    \vspace{0.05cm}
    From closedness of constraints to non-closedness.\par}
\end{frame}

\begin{frame}
  \frametitle{The General Heaviside Composite Optimization Problem\\[-1pt]
    {\normalsize\mdseries (G-HSCOP)}}

  \begingroup
  \setlength{\fboxsep}{0.10cm}
  \setlength{\fboxrule}{0.45pt}
  \color{MediumBlack}
  \noindent\fbox{%
    \begin{minipage}{0.94\textwidth}
      \color{MediumBlue}
      \vspace{-0.05cm}
      \noindent\makebox[\linewidth][c]{\(\displaystyle
        \underset{\mathbf{x}\in P}{\operatorname{maximize}}\quad
        \Phi(\mathbf{x})\triangleq c(\mathbf{x})
        +\sum_{k=1}^{K_0}\psi_{0k}(\mathbf{x})
        \underbrace{\mathds{1}_{[0,\infty)}\!\left(\phi_{0k}(\mathbf{x})\right)}_
          {\text{composite indicator fnc.}}
      \)}\par
      \vspace{0.05cm}
      \noindent\makebox[\linewidth][c]{\(\displaystyle
        \operatorname{subject\ to}\quad
        \mathbf{A}_{i\bullet}\mathbf{x}
        +\sum_{k=1}^{K_i}\psi_{ik}(\mathbf{x})
        \mathds{1}_{[0,\infty)}\!\left(\phi_{ik}(\mathbf{x})\right)
        \geq \eta_i,\qquad i=1,\ldots,I.
      \)}\par
      \vspace{-0.05cm}
    \end{minipage}%
  }
  \endgroup

  \vspace{0.08cm}
  \footnotesize
  \begin{itemize}
    \setlength{\itemsep}{0.02cm}
    \item $P$ is a \textcolor{MediumRed}{friendly} set, $c$ is a
      \textcolor{MediumRed}{friendly} function, e.g., affine or quadratic;
      \textcolor{MediumRed}{friendly} $=$ amenable to well-known treatments,
      and $\mathbf{A}_{i\bullet}$ is the $i$-th row of a matrix
      $\mathbf{A}\in\mathbb{R}^{I\times n}$.
    \item $\psi_{ik}$ and $\phi_{ik}$ are functions whose properties highlight
      the challenges;
    \item most importantly, $\mathds{1}_{[0,\infty)}$ is the ``closed''
      Heaviside function given by
  \end{itemize}

  \vspace{-0.16cm}
  {\color{MediumBlue}\small
  \[
    \mathds{1}_{[0,\infty)}(s)\triangleq
    \begin{cases}
      1, & \text{if }s\in[0,\infty),\\
      0, & \text{otherwise},
    \end{cases}
    \qquad\text{(upper semi-continuous).}
  \]
  }
\end{frame}

\begin{frame}
  \frametitle{The General Heaviside Composite Optimization Problem\\[-1pt]
    {\normalsize\mdseries (G-HSCOP)}}

  \begin{supplementarycallout}{{\small Supplementary visualization: closed vs. open Heaviside thresholds}}
  \small
  \setlength{\abovedisplayskip}{3pt}
  \setlength{\belowdisplayskip}{3pt}
  The Heaviside indicator turns a scalar condition into a binary switch. The
  bracket at $s=0$ determines both activation at equality and the direction
  of semicontinuity.

  \vspace{0.02cm}
  \begin{columns}[T,onlytextwidth]
    \column{0.48\textwidth}
    {\color{MediumBlue}\bfseries Closed threshold: $[0,\infty)$\par}
    {\scriptsize Upper, not lower, semicontinuous; favorable for $\max$.\par}
    \[
      H_{\mathrm{cl}}(s)\triangleq\mathds{1}_{[0,\infty)}(s)
      =\begin{cases}
        0, & s<0,\\
        1, & s\geq 0.
      \end{cases}
    \]
    \vspace{-0.20cm}
    \centering
    \begin{tikzpicture}
      \begin{axis}[
        width=0.96\linewidth,
        height=2.31cm,
        axis lines=middle,
        xmin=-2.1,xmax=2.1,
        ymin=-0.22,ymax=1.32,
        xtick={-2,-1,0,1,2},
        ytick={0,1},
        xlabel={$s$},
        ylabel={$H_{\mathrm{cl}}(s)$},
        label style={font=\scriptsize},
        tick label style={font=\scriptsize},
        axis line style={MediumBlack},
        every axis x label/.style={at={(ticklabel* cs:1)},anchor=west},
        every axis y label/.style={at={(ticklabel* cs:1)},anchor=south},
        clip=false
      ]
        \addplot[MediumBlue,very thick,domain=-2:0] {0};
        \addplot[MediumBlue,very thick,domain=0:2] {1};
        \addplot[MediumBlue,only marks,mark=o,mark size=2.4pt,
          mark options={fill=SupplementTint,very thick}] coordinates {(0,0)};
        \addplot[MediumBlue,only marks,mark=*,mark size=2.4pt]
          coordinates {(0,1)};
      \end{axis}
    \end{tikzpicture}

    \raggedright\footnotesize
    $H_{\mathrm{cl}}(0)=1$; equality activates the term.

    \column{0.48\textwidth}
    {\color{SupplementAccent}\bfseries Open threshold: $(0,\infty)$\par}
    {\scriptsize Lower, not upper, semicontinuous; favorable for $\min$.\par}
    \[
      H_{\mathrm{op}}(s)\triangleq\mathds{1}_{(0,\infty)}(s)
      =\begin{cases}
        0, & s\leq 0,\\
        1, & s>0.
      \end{cases}
    \]
    \vspace{-0.20cm}
    \centering
    \begin{tikzpicture}
      \begin{axis}[
        width=0.96\linewidth,
        height=2.31cm,
        axis lines=middle,
        xmin=-2.1,xmax=2.1,
        ymin=-0.22,ymax=1.32,
        xtick={-2,-1,0,1,2},
        ytick={0,1},
        xlabel={$s$},
        ylabel={$H_{\mathrm{op}}(s)$},
        label style={font=\scriptsize},
        tick label style={font=\scriptsize},
        axis line style={MediumBlack},
        every axis x label/.style={at={(ticklabel* cs:1)},anchor=west},
        every axis y label/.style={at={(ticklabel* cs:1)},anchor=south},
        clip=false
      ]
        \addplot[SupplementAccent,very thick,domain=-2:0] {0};
        \addplot[SupplementAccent,very thick,domain=0:2] {1};
        \addplot[SupplementAccent,only marks,mark=*,mark size=2.4pt]
          coordinates {(0,0)};
        \addplot[SupplementAccent,only marks,mark=o,mark size=2.4pt,
          mark options={fill=SupplementTint,very thick}] coordinates {(0,1)};
      \end{axis}
    \end{tikzpicture}

    \raggedright\footnotesize
    $H_{\mathrm{op}}(0)=0$; strict inequality is required. In particular,
    $|t|_0=H_{\mathrm{op}}(|t|)$.
  \end{columns}

  \vspace{-0.02cm}
  {\color{MediumBlue}\footnotesize
    Thus $\psi(\mathbf{x})H_{\mathrm{cl}}(\phi(\mathbf{x}))$ equals
    $\psi(\mathbf{x})$ when $\phi(\mathbf{x})\geq0$, and equals $0$
    otherwise.}
  \end{supplementarycallout}
\end{frame}

\begin{frame}{Conditional Functions and Sparsity}
  \small
  \setlength{\abovedisplayskip}{4pt}
  \setlength{\belowdisplayskip}{4pt}
  \setlength{\abovedisplayshortskip}{3pt}
  \setlength{\belowdisplayshortskip}{3pt}
  The product
  $f(\mathbf{x})\mathds{1}_{[0,\infty)}\!\left(g(\mathbf{x})\right)$
  is a conditional function; i.e., the primary function $f(\mathbf{x})$ is
  activated only if the auxiliary condition $g(\mathbf{x})\geq 0$ holds.

  \vspace{0.12cm}
  {\color{MediumBlue}An immediate example: the sparsity (counting) function}
  \vspace{-0.04cm}
  \[
    |t|_0\triangleq
    \begin{cases}
      1, & \text{if }t\neq 0,\\
      0, & \text{if }t=0,
    \end{cases}
    =\mathds{1}_{(0,\infty)}(|t|)
    =1-\mathds{1}_{[0,\infty)}(-|t|),
  \]
  \vspace{-0.08cm}

  illustrates that the function inside the Heaviside function may not be
  differentiable. This function extends to any step function with a finite
  number of (discontinuous) jumps.

  \vspace{0.14cm}
  An \textcolor{MediumRed}{affine sparsity} constraint
  \vspace{-0.05cm}
  \[
    \sum_{j=1}^{n}a_j|x_j|_0\leq b,
    \qquad \text{with $a_j$ of mixed signs},
  \]
  \vspace{-0.08cm}

  is a non-trivial extension of the cardinality constraint
  $\sum_{j=1}^{n}|x_j|_0\leq K$, whose mixing coefficients are all positive.

  \begin{supplementarynote}
    Positive coefficients preserve lower semicontinuity. With mixed-sign
    $a_j$, upward and downward jumps can coexist, so
    $\sum_{j=1}^{n}a_j|x_j|_0$ may be neither lower nor upper
    semicontinuous.
  \end{supplementarynote}
\end{frame}

\begin{frame}{References}
  \footnotesize
  \setlength{\parskip}{0pt}

  {\color{MediumRed}\bfseries Methodology}
  \vspace{0.04cm}
  \begin{itemize}
    \setlength{\itemsep}{0.08cm}
    \item Y.~\textsc{Cui}, J.~\textsc{Liu}, \textsc{and} J.S.~\textsc{Pang}.
      On the minimization of piecewise functions: Pseudo stationarity.
      \textit{Journal of Convex Analysis} 30(3): 793--834 (2023).
    \item S.~\textsc{Han}, Y.~\textsc{Cui}, \textsc{and} J.S.~\textsc{Pang}.
      Analysis of a class of minimization problems lacking lower
      semicontinuity. \textit{Mathematics of Operations Research}
      50(3): 2175--2198 (2025-8).
    \item K.~\textsc{Zheng}, J.~\textsc{Liu}, R.~\textsc{Wang}, \textsc{and}
      J.S.~\textsc{Pang}. Solving constrained affine Heaviside composite
      optimization problems by a progressive IP approach.
      Manuscript (April 2026).
  \end{itemize}

  \vspace{0.20cm}
  {\color{MediumRed}\bfseries Applications}
  \vspace{0.04cm}
  \begin{itemize}
    \setlength{\itemsep}{0.08cm}
    \item Z.~\textsc{Qi}, Y.~\textsc{Cui}, Y.~\textsc{Liu}, \textsc{and}
      J.S.~\textsc{Pang}. Estimation of individualized decision rules based on
      optimized covariate-dependent equivalent of random outcomes.
      \textit{SIAM Journal on Optimization} 29(3): 2337--2362 (2019).
    \item Y.~\textsc{Fang}, J.~\textsc{Liu}, \textsc{and} J.S.~\textsc{Pang}.
      Treatment learning with Gini constraints by Heaviside composite
      optimization and a progressive method.
      \textit{Computational Optimization and Applications}
      92(2): 471--513 (2025-11).
  \end{itemize}
\end{frame}

\begin{frame}{Distinguished (and challenging) features:}
  \small
  \setlength{\parskip}{0pt}
  \begin{itemize}
    \setlength{\itemsep}{0.14cm}
    \item \textcolor{MediumRed}{Discontinuity} of the Heaviside function and
      \textcolor{MediumRed}{nondifferentiability} of the other functions:
      \begin{itemize}
        \setlength{\itemsep}{0.04cm}
        \item[---] lack of semicontinuity in the objective function
        \item[---] non-closedness of the feasible set
        \item[---] does convexity/concavity have a role to play?
      \end{itemize}

    \item \textcolor{MediumRed}{Composition} of the highlighted features:
      \begin{itemize}
        \setlength{\itemsep}{0.04cm}
        \item[---] what is a reasonable concept of a
          \textcolor{MediumRed}{solution}?
        \item[---] algorithmically, how to compute such a solution?
      \end{itemize}

    \item \textcolor{MediumRed}{Fusion} of continuous variables and functions
      with discrete structure of the Heaviside function; thus natural to
      consider
      \begin{itemize}
        \setlength{\itemsep}{0.04cm}
        \item[---] a \textcolor{MediumRed}{marriage} between the continuous
          and discrete world
        \item[---] how to facilitate the marriage? will it be fruitful?
      \end{itemize}
  \end{itemize}
\end{frame}

\section{Chapter II: Prerequisites, a second course in optimization}

\begin{frame}[plain]
  \centering
  \vfill
  {\usebeamercolor[fg]{structure}\usebeamerfont{title}Chapter II\par}
  \vspace{0.34cm}
  {\usebeamercolor[fg]{structure}\usebeamerfont{subtitle}\mdseries\large
    Prerequisites, a second course in optimization\par}
  \vfill
\end{frame}

% Original slide 11.  The rest of Chapter II (slides 10, 14--20) is still
% awaiting photographs.
\begin{frame}{Prerequisites, a second course in optimization}
  \small
  \setlength{\parskip}{0pt}
  \begin{itemize}
    \setlength{\itemsep}{0.16cm}
    \item \textcolor{MediumRed}{A graduate course in linear, differentiable,
      and integer programming}
    \item \textbf{Bouligand differentiability}: broad and most basic for
      nondifferentiable functions
    \item \textbf{Stationarity conditions} of nondifferentiable programs
    \item \textbf{Piecewise affine functions}: definition and properties
    \item \textbf{Difference-of-convex programs}: criticality vs.\
      stationarity; the difference-of-convex algorithm and its refinement
    \item \textbf{Complementarity constraints}: modeling breadth and
      connections to binary programming
    \item \textbf{Fractional and composite programs}: Dinkelbach's algorithm
      and its extension.
  \end{itemize}

  \begin{supplementarynote}
    $0\leq a\perp b\geq0$ is an or-gate --- $a_i=0$ \emph{or} $b_i=0$ ---
    i.e.\ if-then logic written in continuous variables.
  \end{supplementarynote}
\end{frame}

% Original slide 12.
\begin{frame}{Bouligand differentiability}
  \small
  \setlength{\abovedisplayskip}{5pt}
  \setlength{\belowdisplayskip}{5pt}

  A function $f:\mathcal{O}\subseteq\mathbb{R}^n\to\mathbb{R}$, where
  $\mathcal{O}$ is an open set in $\mathbb{R}^n$ is
  {\color{MediumBlue}Bouligand differentiable} at
  $\bar{\mathbf{x}}\in\mathcal{O}$ if it is {\color{MediumBlue}locally
  Lipschitz continuous} near $\bar{\mathbf{x}}$ and directionally
  differentiable there with the (finite) directional derivative along the
  direction $\mathbf{v}\in\mathbb{R}^n$ defined by:
  \[
    \fp (\bar{\mathbf{x}};\mathbf{v})\triangleq\lim_{\tau\downarrow0}
    \frac{f(\bar{\mathbf{x}}+\tau\mathbf{v})-f(\bar{\mathbf{x}})}{\tau}.
  \]

  As a function of the direction, $\fp (\bar{\mathbf{x}};\bullet)$ is
  positively homogeneous and globally Lipschitz continuous on $\mathbb{R}^n$
  with the same Lipschitz modulus as that of $f$ near $\bar{\mathbf{x}}$.

  \vspace{0.06cm}
  Calculus rules hold for the directional derivative; these are the sum,
  multiplication, division, and composition rules; in particular,
  \[
    (f\circ g)'(\bar{\mathbf{x}};\mathbf{v})
    =\fp \bigl(g(\bar{\mathbf{x}});g'(\bar{\mathbf{x}};\mathbf{v})\bigr),
  \]

  \textbf{Remark:} the partial differentiation rule may fail;
  \[
    \fp (\bar{\mathbf{x}};\mathbf{v})\neq\sum_{i=1}^{n}
    \fp (\bar{\mathbf{x}};v_i\mathbf{e}^i),\qquad
    \text{where }\ \mathbf{e}^i=i\text{-th coordinate vector}
  \]

  \begin{supplementarynote}
    Convex functions are directionally differentiable.
  \end{supplementarynote}
\end{frame}

% Original slide 13.
\begin{frame}{Piecewise affine functions}
  \small
  \setlength{\abovedisplayskip}{9pt}
  \setlength{\belowdisplayskip}{9pt}

  A {\color{MediumBlue}continuous} function $f:\mathbb{R}^n\to\mathbb{R}^m$
  is {\color{MediumBlue}piecewise affine} if there exists a partition of
  $\mathbb{R}^n$ into a finite number of polyhedra, on each of which, $f$ is
  an affine function.

  \vspace{0.12cm}
  A piecewise affine function is Bouligand differentiable everywhere.
  \textcolor{MediumRed}{Most interested in $m=1$.}

  \vspace{0.12cm}
  A \textcolor{MediumRed}{continuous} function
  $f:\mathbb{R}^n\to\mathbb{R}^m$ is piecewise affine (PA) if and only if
  \[
    f(\mathbf{x})=\underbrace{\max_{1\leq k\leq K}
      \bigl[(\mathbf{a}^k)^{\top}\mathbf{x}+\alpha_k\bigr]}_
      {\text{convex in }\mathbf{x}}
    -\underbrace{\max_{1\leq\ell\leq L}
      \bigl[(\mathbf{b}^{\ell})^{\top}\mathbf{x}+\beta_{\ell}\bigr]}_
      {\text{convex in }\mathbf{x}},\qquad\mathbf{x}\in\mathbb{R}^n
  \]

  for some constant vectors $\mathbf{a}^k$ and $\mathbf{b}^{\ell}$ and
  scalars $\alpha_k$ and $\beta_{\ell}$. Hence, if $f$ is PA, then for every
  $\bar{\mathbf{x}}\in\mathbb{R}^n$,
  \[
    f(\mathbf{x})=f(\bar{\mathbf{x}})
    +\fp (\bar{\mathbf{x}};\mathbf{x}-\bar{\mathbf{x}}),\qquad
    \forall\,\mathbf{x}\text{ sufficiently near }\bar{\mathbf{x}},
  \]
\end{frame}

\section{Chapter III: Paradigm changes -- Sources of discontinuity}

\begin{frame}[plain]
  \centering
  \vfill
  {\usebeamercolor[fg]{structure}\usebeamerfont{title}Chapter III\par}
  \vspace{0.34cm}
  {\usebeamercolor[fg]{structure}\usebeamerfont{subtitle}\mdseries\large
    Paradigm changes: Sources of discontinuity\par}
  \vfill
\end{frame}

\begin{frame}
  \frametitle{Modeling breadth
    {\normalsize\mdseries (OR, statistics, ML, and policy learning)}}

  \small
  \setlength{\parskip}{0pt}
  \begin{itemize}
    \setlength{\itemsep}{0.13cm}
    \item {\color{MediumBlue}Set-up costs}--- classical in operations research,
      today's term is sparsity
    \item {\color{MediumBlue}Structured variable/constraint selection};
      on-off constraints
    \item {\color{MediumBlue}Conditional stochastic functionals} ---
      probabilities and expectations
    \item {\color{MediumBlue}Statistical models} --- conformal regression;
      change-point detection; robust trimming; and more
    \item {\color{MediumBlue}Classification} --- score-based binary
      classification; multiclass classification; advanced measures;
      decision trees
    \item {\color{MediumBlue}Clustering: refreshed and extended} ---
      mixed-norm: extending traditional K-means and K-clustering;
      bicriteria (errors versus similarities)
    \item {\color{MediumBlue}Univariate and vector quantization}
    \item {\color{MediumBlue}Policy learning and causal inference} ---
      piecewise affine policy for stochastic optimization; causal inference
      for individualized treatment learning; policy learning with optimal
      weight clipping
  \end{itemize}
\end{frame}

% A topic divider inside Chapter II -- deliberately not a \section, so it
% sits one step below the chapter dividers in the visual hierarchy.
\begin{frame}[plain]
  \centering
  \vfill
  {\usebeamercolor[fg]{structure}\Large\bfseries A host of instances\par}
  \vspace{0.26cm}
  {\usebeamercolor[fg]{structure}\normalsize\mdseries
    from classical to modern to future\par}
  \vfill
\end{frame}

\begin{frame}{Hierarchical variable selection}
  \small
  \setlength{\abovedisplayskip}{4pt}
  \setlength{\belowdisplayskip}{4pt}
  \setlength{\abovedisplayshortskip}{3pt}
  \setlength{\belowdisplayshortskip}{3pt}

  The statistical model with interaction terms: hypothesizing a bilinear
  relation among the input features
  $\boldsymbol{\xi}=(\xi_i)_{i=1}^{d}$
  \[
    m(\mathbf{x};\boldsymbol{\xi})
    =\sum_{i=1}^{d}x_i^{(1)}\xi_i
    +\sum_{1\leq i<j\leq d}x_{ij}^{(2)}\xi_i\xi_j,
  \]
  with the parameter $\mathbf{x}$ consisting of the unknown model
  coefficients $x_i^{(1)}$ and $x_{ij}^{(2)}$ that obey

  \begin{itemize}
    \setlength{\itemsep}{0.08cm}
    \item {\color{MediumBlue}strong hierarchy}: an interaction term can be
      selected only if both of the linear terms are selected, i.e.,
      \[
        \bigl|x_{ij}^{(2)}\bigr|_0\leq
        \min\bigl\{\bigl|x_i^{(1)}\bigr|_0,\bigl|x_j^{(1)}\bigr|_0\bigr\};
      \]
    \item {\color{MediumBlue}weak hierarchy}: an interaction term can be
      selected only if at least one of the corresponding linear terms is
      selected,
      \[
        \bigl|x_{ij}^{(2)}\bigr|_0\leq
        \bigl|x_i^{(1)}\bigr|_0+\bigl|x_j^{(1)}\bigr|_0;
      \]
  \end{itemize}

  both leading to \textcolor{MediumRed}{mixed-sign} combinations of
  Heaviside functions.
\end{frame}

\begin{frame}{Conditional stochastic functionals}
  \small
  \setlength{\abovedisplayskip}{5pt}
  \setlength{\belowdisplayskip}{5pt}
  \setlength{\abovedisplayshortskip}{4pt}
  \setlength{\belowdisplayshortskip}{4pt}

  \begin{itemize}
    \setlength{\itemsep}{0.12cm}
    \item Expression of a probability function as an expectation:
      \[
        \mathbb{P}\bigl[f(\mathbf{x},\boldsymbol{\xi})\geq 0\bigr]
        =\mathbb{E}\,\mathds{1}_{[0,\infty)}
          \bigl(f(\mathbf{x},\boldsymbol{\xi})\bigr)
        \xrightarrow[\text{approximation (SAA)}]{\text{sample average}}
        \frac{1}{N}\sum_{s=1}^{N}\mathds{1}_{[0,\infty)}
          \bigl(f(\mathbf{x},\boldsymbol{\xi}^s)\bigr);
      \]
    \item and by extension, expression of a
      {\color{MediumBlue}decision-dependent conditional expectation
      function}:
      \[
        \mathbb{E}\bigl[\psi(\mathbf{x},\boldsymbol{\xi})\bigm|
          f(\mathbf{x},\boldsymbol{\xi})\geq 0\bigr]
        \triangleq
        \underbrace{\frac{\mathbb{E}\bigl[\psi(\mathbf{x},\boldsymbol{\xi})
            \mathds{1}_{[0,\infty)}
            (f(\mathbf{x},\boldsymbol{\xi}))\bigr]}
          {\mathbb{P}\bigl[f(\mathbf{x},\boldsymbol{\xi})\geq 0\bigr]}}_
          {\text{fractional Heaviside function}}.
      \]
  \end{itemize}

  \vspace{0.10cm}
  Note the product
  \vspace{-0.02cm}
  \[
    \psi(\mathbf{x},\boldsymbol{\xi})\,
    \mathds{1}_{[0,\infty)}\bigl(f(\mathbf{x},\boldsymbol{\xi})\bigr)
  \]

  raising the possibility of losing semicontinuity (in the decision variable
  $\mathbf{x}$) if the sign of the multiplier function
  $\psi(\bullet,\boldsymbol{\xi})$ is not consistent with the optimization.
\end{frame}

\begin{frame}{Mixed-sign combinations and piecewise inner functions}
  \small
  \setlength{\abovedisplayskip}{4pt}
  \setlength{\belowdisplayskip}{4pt}
  \setlength{\abovedisplayshortskip}{3pt}
  \setlength{\belowdisplayshortskip}{3pt}

  \begin{itemize}
    \setlength{\itemsep}{0.16cm}
    \item First-order stochastic dominance:
      \[
        \mathbb{P}\bigl[f_1(\mathbf{x},\boldsymbol{\xi})\leq c_1\bigr]
        \leq
        \mathbb{P}\bigl[f_2(\mathbf{x},\boldsymbol{\xi})\leq c_2\bigr]
      \]
    \item Conjunctive probability constraint:
      \begin{align*}
        b\;&\geq\;\mathbb{P}\bigl[\psi_1(\mathbf{x},\boldsymbol{\xi})\geq 0
          \ \text{ and }\ \psi_2(\mathbf{x},\boldsymbol{\xi})<0\bigr]\\
        &=\;\mathbb{P}\bigl[\psi_1(\mathbf{x},\boldsymbol{\xi})\geq 0\bigr]
          -\mathbb{P}\bigl[\min\bigl(\psi_1(\mathbf{x},\boldsymbol{\xi}),
          \psi_2(\mathbf{x},\boldsymbol{\xi})\bigr)\geq 0\bigr]
      \end{align*}
    \item Conditional probability constraint:
      \vspace{-0.06cm}
      \[
        b\geq\mathbb{P}\bigl[\psi_1(\mathbf{x},\boldsymbol{\xi})\geq 0
          \bigm|\psi_2(\mathbf{x},\boldsymbol{\xi})\geq 0\bigr]
      \]
      \vspace{-0.42cm}
      \[
        \Updownarrow
      \]
      \vspace{-0.42cm}
      \[
        \mathbb{P}\bigl[\min\bigl(\psi_1(\mathbf{x},\boldsymbol{\xi}),
          \psi_2(\mathbf{x},\boldsymbol{\xi})\bigr)\geq 0\bigr]
        -b\,\mathbb{P}\bigl[\psi_2(\mathbf{x},\boldsymbol{\xi})\geq 0\bigr]
        \leq 0.
      \]
  \end{itemize}
\end{frame}

\begin{frame}
  \frametitle{Some statistical models
    {\normalsize\mdseries (communicated by Dr.~Yao Xie)}}

  \small
  \setlength{\abovedisplayskip}{4pt}
  \setlength{\belowdisplayskip}{4pt}

  \begin{itemize}
    \item \textbf{Change-point detection with sparsity.} Let
      $\{\mathbf{x}_t\}_{t=1}^{T}$ be a sequence with
      $\mathbf{x}_t\in\mathbb{R}^p$ and multiple unknown change point
      $\{\tau_i\}_{i=1}^{I}$ with mean shifts
      $\{\boldsymbol{\Delta}_i\}_{i=1}^{I}$ to be determined. Let
      $\boldsymbol{\mu}\in\mathbb{R}^p$ be the baseline mean vector which is
      unknown. This problem is formulated as follows:
  \end{itemize}

  \[
    \begin{aligned}
      \underset{\boldsymbol{\mu},\{\boldsymbol{\Delta}_i\},\{\tau_i\}}
        {\operatorname{minimize}}\quad
      &\sum_{t=1}^{T}\mathds{1}_{(0,\infty)}(\tau_1-t)
        \bigl\|\mathbf{x}_t-\boldsymbol{\mu}\bigr\|_2^2
        +\sum_{i=1}^{I}\bigl\|\boldsymbol{\Delta}_i\bigr\|_0\\[2pt]
      &+\sum_{i=1}^{I}\sum_{t=1}^{T}
        \boxed{\underbrace{\mathds{1}_{(0,\infty)}(t-\tau_i)\,
          \mathds{1}_{[0,\infty)}(\tau_{i+1}-t)}_
          {\text{expressing }\tau_i<t\leq\tau_{i+1}}}\;
        \bigl\|\mathbf{x}_t-\boldsymbol{\mu}-\boldsymbol{\Delta}_i\bigr\|_2^2
        \\[2pt]
      \operatorname{subject\ to}\quad
      &\tau_1\leq\tau_2\leq\cdots\leq\tau_I.
    \end{aligned}
  \]

  \textcolor{MediumRed}{Note} the product of an open and a closed Heaviside
  function.

  \begin{supplementarynote}
    Matched brackets let two factors collapse into one,
    $\mathds{1}_{[0,\infty)}(a)\,\mathds{1}_{[0,\infty)}(b)
     =\mathds{1}_{[0,\infty)}\bigl(\min(a,b)\bigr)$, and likewise for
    $\mathds{1}_{(0,\infty)}$. The mixed pair above does not, since
    $\{a>0,\,b\geq0\}$ is neither open nor closed. Under
    $\tau_i\leq\tau_{i+1}$ it is instead the mixed-sign difference
    $\mathds{1}_{[0,\infty)}(\tau_{i+1}-t)
     -\mathds{1}_{[0,\infty)}(\tau_i-t)$, with affine inner functions.
  \end{supplementarynote}
\end{frame}

\begin{frame}
  \frametitle{Some statistical models
    {\normalsize\mdseries (cont.)}}

  \small
  \setlength{\abovedisplayskip}{6pt}
  \setlength{\belowdisplayskip}{6pt}

  \begin{itemize}
    \item \textbf{Trimmed least-squares for robust statistical learning.}
  \end{itemize}

  \[
    \begin{aligned}
      \underset{f\in\mathcal{F};\,\gamma\geq0}{\operatorname{minimize}}\quad
      &\sum_{i=1}^{n}\underbrace{\mathds{1}_{[0,\infty)}\bigl(
         \underbrace{\gamma}_{\mathclap{\text{trimming threshold}}}
         -\|\mathbf{y}_i-f(\mathbf{x}_i)\|\bigr)
         \bigl\|\mathbf{y}_i-f(\mathbf{x}_i)\bigr\|_2^2}_
         {\text{a conditional error}}\\[16pt]
      \operatorname{subject\ to}\quad
      &\underbrace{\frac{1}{n}\sum_{i=1}^{n}\mathds{1}_{[0,\infty)}\bigl(
         \gamma-\|\mathbf{y}_i-f(\mathbf{x}_i)\|\bigr)\geq 1-\beta}_
         {\text{bounding the trimming ratio}}.
    \end{aligned}
  \]
\end{frame}

\begin{frame}
  \frametitle{Clustering problems:
    {\normalsize\mdseries Classical and extended.}}

  \footnotesize
  \setlength{\abovedisplayskip}{3pt}
  \setlength{\belowdisplayskip}{3pt}
  \setlength{\abovedisplayshortskip}{2pt}
  \setlength{\belowdisplayshortskip}{2pt}

  Let $\Xi\triangleq\{\boldsymbol{\xi}^s\}_{s=1}^{N}$ be a given set of data
  points in $\mathbb{R}^p$. Let
  $\mathbf{c}\triangleq\{\mathbf{c}^k\}_{k=1}^{K}$ be an unknown set of
  cluster centers in $\mathbb{R}^p$. The criterion for a data point
  $\boldsymbol{\xi}^s$ to belong to cluster $k$, written as $s\in k$, is
  \[
    \begin{aligned}
      &\bigl\|\boldsymbol{\xi}^s-\mathbf{c}^k\bigr\|\leq
        \min_{k'\neq k}\bigl\|\boldsymbol{\xi}^s-\mathbf{c}^{k'}\bigr\|\\
      \Leftrightarrow\ \ &h_s(\mathbf{c})\triangleq
        \min_{k'\neq k}\bigl\|\boldsymbol{\xi}^s-\mathbf{c}^{k'}\bigr\|
        -\bigl\|\boldsymbol{\xi}^s-\mathbf{c}^k\bigr\|\geq 0.
    \end{aligned}
  \]

  Define the simple indicator
  \[
    z_{sk}\triangleq\mathds{1}(s\in k)\triangleq
    \mathds{1}_{[0,\infty)}\bigl\{h_s(\mathbf{c})\bigr\}
    =\begin{cases}
      1 & \text{if }s\in k\\
      0 & \text{if }s\notin k
    \end{cases}
  \]

  \vspace{0.06cm}
  {\color{MediumBlue}The classical K-means $(p=2)$ and K-medians $(p=1)$
  problems:}
  \[
    \begin{aligned}
      \underset{\mathbf{c};\,\mathbf{z}}{\operatorname{minimize}}\quad
      &\frac{1}{N}\sum_{s=1}^{N}\Biggl[\,\sum_{k=1}^{K}z_{sk}
        \bigl\|\boldsymbol{\xi}^s-\mathbf{c}^k\bigr\|_p^p\Biggr]\\
      \operatorname{subject\ to}\quad
      &\sum_{k=1}^{K}z_{sk}=1,\quad\forall\,s=1,\ldots,N
        &&\begin{array}{@{}l@{}}
            \text{\scriptsize every sample belongs to}\\[-2pt]
            \text{\scriptsize at least one cluster}
          \end{array}\\
      \text{and}\quad
      &z_{sk}\in\{0,1\},\quad\forall\,(s,k)\in[N]\times[K].
    \end{aligned}
  \]
\end{frame}

\begin{frame}
  \frametitle{Clustering problems:
    {\normalsize\mdseries Classical and extended.} (cont.)}

  \footnotesize
  \setlength{\abovedisplayskip}{2pt}
  \setlength{\belowdisplayskip}{2pt}

  Unifying the above two problems is the mixed $\ell_{p,q}$-clustering
  problem with arbitrary pairs of integers $p,q\geq 1$:
  \vspace{-0.16cm}
  \[
    \begin{aligned}
      \underset{\mathbf{c},\mathbf{z}}{\operatorname{minimize}}\quad
      &\underbrace{\frac{1}{N}\sum_{s=1}^{N}\Biggl[\,\sum_{k=1}^{K}z_{sk}
        \bigl\|\boldsymbol{\xi}^s-\mathbf{c}^k\bigr\|_p^p\Biggr]}_
        {\text{bi-convex; denoted WCD}_p(\mathbf{c},\mathbf{z})}
      &&{\color{MediumBlue}\ell_p\text{-norm in objective}}\\[1pt]
      \operatorname{subject\ to}\quad
      &\sum_{k=1}^{K}z_{sk}=1,\quad\forall\,s\in[N]\triangleq
        \{1,\ldots,N\}\\
      \text{and}\quad&\text{for all }(s,k)\in[N]\times[K]:\\
      &\underbrace{z_{sk}\leq\mathds{1}_{[0,\infty)}\Bigl(
        \min_{k'\neq k}\bigl\|\boldsymbol{\xi}^s-\mathbf{c}^{k'}\bigr\|_q
        -\bigl\|\boldsymbol{\xi}^s-\mathbf{c}^{k}\bigr\|_q\Bigr)}_
        {\textcolor{MediumRed}{\text{Heaviside constraint with }
          \ell_q\text{-norm}}}\\[1pt]
      &z_{sk}\in\{0,1\},\quad\text{(binary variables)}
    \end{aligned}
  \]

  \textcolor{MediumRed}{Note:} The Heaviside constraint is not needed in the
  classical K-means and K-medians problems where $p=q$.

  \begin{supplementarynote}
    When $p=q$ the objective alone decides membership. Writing the indicator
    out separates the cost ($\ell_p$) from the rule ($\ell_q$), making the
    rule a free modeling choice.
  \end{supplementarynote}
\end{frame}

\end{document}
